\lysh{38857}{4}

\begin{alist}
\item
\begin{rlist}
\item Αρχικά αγοράζει μέλι ένας πελάτης και στη συνέχεια αυτός συστήνει το προϊόν και αγοράζουν μέλι $2$ καινούριοι πελάτες και αυτοί με τη σειρά τους συστήνουν σε άλλους $2$ ο καθένας, άρα $4$ συνολικά κ.ο.κ.
Άρα έχουμε τη γεωμετρική πρόοδο: $1, 2, 4, \dots$ με πρώτο όρο $\alpha_1 = 1$ και λόγο $\lambda = 2$.
\item Για να βρούμε τον $\nu$-οστό όρο της γεωμετρικής προόδου, θα χρησιμοποιήσουμε τον τύπο:
\[\alpha_\nu = \alpha_1 \cdot \lambda^{\nu-1} = 1 \cdot 2^{\nu-1} = 2^{\nu-1}.\]
Για να υπολογίσουμε τον 7ο όρο θα αντικαταστήσουμε όπου $\nu = 7$ και έχουμε:
$\alpha_7 = 2^{7-1} = 2^6 = 64$, άρα $64$ καινούριοι πελάτες θα αγοράσουν μέλι στην 7η φάση της διαδικασίας.
\item Μετά από τη $\nu$-οστή φάση αυτής της διαδικασίας θα έχουν πουληθεί συνολικά
$S_\nu = \alpha_1 \cdot \dfrac{\lambda^\nu - 1}{\lambda - 1} = 1 \cdot \dfrac{2^\nu - 1}{2 - 1} = 2^\nu - 1$ βάζα μέλι.

Άρα, $S_7 = 2^7 - 1 = 127$. Άρα, αν πετύχει η διαδικασία προώθησης που σκέφτηκε η Άννα, μετά από 7 φάσεις που θα έχουν πραγματοποιηθεί με επιτυχία, θα έχει πουλήσει 127 βάζα μελιού συνολικά.

\textit{Εναλλακτικός τρόπος:} Υπολογίζοντας ξεχωριστά και αθροίζοντας τους 7 πρώτους όρους: $\alpha_1 + \alpha_2 + \alpha_3 + \alpha_4 + \alpha_5 + \alpha_6 + \alpha_7 = 1 + 2 + 4 + 8 + 16 + 32 + 64 = 127$.
\end{rlist}

\item Αφού θέλει να διεξαχθούν το πολύ 9 φάσεις αυτής της διαδικασίας, συνολικά οι πελάτες που θα αγοράσουν μέλι με αυτή τη διαδικασία, εφόσον πετύχει σε όλες τις φάσεις της, θα είναι το πολύ: $S_9 = 2^9 - 1 = 512 - 1 = 511$.
Η 9η φάση λοιπόν δεν πρέπει να ολοκληρωθεί, καθώς τότε θα υπερβεί το πλήθος των 500 πελατών.
\end{alist}
